% Generated draft: requires mathematical and source review.
\documentclass[11pt]{article}
\usepackage[margin=25mm]{geometry}
\usepackage{fontspec}
\setmainfont{Latin Modern Roman}
\usepackage{amsmath,amssymb,mathtools,unicode-math}
\setmathfont{Latin Modern Math}
\usepackage{enumitem,needspace,xurl,hyperref}
\setlength{\parindent}{0pt}
\setlength{\parskip}{5pt}
\setlength{\emergencystretch}{3em}
\newcommand{\N}{\mathbb{N}}
\newcommand{\Z}{\mathbb{Z}}
\newcommand{\Q}{\mathbb{Q}}
\newcommand{\R}{\mathbb{R}}
\newcommand{\C}{\mathbb{C}}
\newcommand{\F}{\mathbb{F}}
\newcommand{\A}{\mathbb{A}}
\newcommand{\PP}{\mathbb{P}}
\newcommand{\E}{\mathbb{E}}
\newcommand{\eps}{\varepsilon}
\newcommand{\dd}{\,\mathrm{d}}
\begin{document}
\section*{UnsolvedMath ID: 30005582}
\textbf{Source problem code:} catalogue entry
\subsection{Definitions and mathematical statement}
Let \(n\geq 1\) and \(T>0\). For \(i\in\{1,2\}\), let \(M_i^n\) be a smooth connected manifold, possibly incomplete, and let \(g_i(t)\), \(0<t<T\), be a smooth Ricci flow, so that
\[
\partial_t g_i(t)=-2\operatorname{Ric}(g_i(t)).
\]
The source additionally requires both flows to satisfy its conditions \((1)\) and \((2)\). Those conditions are absent from the supplied record; they remain unresolved hypotheses of this draft and must not be omitted from a completed formulation.

Write \(d_{g_i(t)}\) for the distances associated with the evolving metrics. On neighborhoods \(U_{p_i,i}\subset M_i\), suppose that there are limiting distances \(d_{0,i}\) such that, in the source's unspecified sense of local convergence,
\[
\lim_{t\downarrow0}d_{g_i(t)}=d_{0,i}\quad\text{on }U_{p_i,i}.
\]
Suppose that a homeomorphism
\[
\psi_0:U_{p_1,1}\longrightarrow U_{p_2,2},\qquad
\psi_0(U_{p_1,1})=U_{p_2,2},
\]
identifies these limiting metric spaces isometrically, meaning
\[
d_{0,2}\bigl(\psi_0(x),\psi_0(y)\bigr)=d_{0,1}(x,y)
\quad\text{for all }x,y\in U_{p_1,1}.
\]
Denote the identified local limiting distance by \(d_0\). The record does not specify whether the distances on these neighborhoods are restrictions of ambient distances or intrinsic neighborhood distances.

\textbf{Problem (conditional on the missing source hypotheses).} Determine what further regularity assumptions on \(d_0\) ensure that the two flows admit a suitable common gauge in which they can be effectively compared and remain close for times tending to zero. A precise theorem must specify the gauge, its spatial domain and time interval, the comparison topology or norm, and the quantitative meaning of remaining close. These are requested outcomes of the source problem, rather than supplied definitions. In particular, equality of the flows or uniqueness is not an explicitly stated conclusion of the extracted question.

The source's fundamental regularity assumption is the following condition, labeled \((3)\): for every \(p\in M\) and every \(x\in U_p\), every tangent cone at \(x\) of \((U_p,d_0)\) is isometric to
\[
(\R^n,d_\delta),\qquad d_\delta(u,v)=|u-v|,
\]
where \(\delta\) is the Euclidean metric. In the usual metric interpretation, a tangent cone is a pointed Gromov--Hausdorff limit of rescalings \((U_p,r_j^{-1}d_0,x)\) along a sequence \(r_j\downarrow0\), when such a limit exists. The record does not supply the hypotheses ensuring existence of these limits. It describes \((3)\) as an infinitesimal absence of singular points.

The source further states that if \((M,d_0)\) is the local limit of a non-collapsing sequence of complete smooth Riemannian manifolds with bounded curvature and Ricci curvature uniformly bounded from below, then \((3)\) can be strengthened to the following uniform Reifenberg condition, labeled \((4)\). For every \(p\in M\) and every \(\varepsilon>0\), there exist \(r>0\) and a neighborhood \(U_p\Subset M\) such that
\[
 d_{\mathrm{GH}}\!\left(
 \bigl(B_{d_0}(x,s),s^{-1}d_0\bigr),
 \bigl(B_{d_\delta}(0,1),d_\delta\bigr)
 \right)<\varepsilon
\]
for every \(0<s<r\) and every \(x\in U_p\) satisfying \(B_{d_0}(x,s)\Subset U_p\). Here \(\Subset\) denotes compact containment and \(d_{\mathrm{GH}}\) denotes Gromov--Hausdorff distance. This rewrites the source's ball of radius one for the rescaled distance \(s^{-1}d_0\) as a ball of radius \(s\) with its distance rescaled. The source does not specify open versus closed balls, the precise local-limit topology, the non-collapsing convention, or the uniformity intended by bounded curvature. Consequently this strengthening is retained as a qualified source assertion requiring review, rather than asserted here as an independently verified theorem.

The source also emphasizes that \((3)\) need not imply that \(d_0\) is induced by a H\"older continuous Riemannian metric, let alone a Lipschitz Riemannian metric. It cites its reference \([5,\text{ Theorem }1.2]\) for a counterexample, but the corresponding bibliographic entry is not supplied.
\subsection{Short English statement}
Determine which regularity assumptions on a common local metric-space limit allow two smooth, connected, possibly incomplete Ricci flows to be compared in a suitable gauge and to remain close as \(t\downarrow0\). The source emphasizes Euclidean tangent cones and a qualified uniform Reifenberg condition. Its hypotheses \((1)\) and \((2)\), convergence conventions, gauge, and closeness estimates are missing or unspecified, so this remains a conditional statement-only draft requiring mathematical review.
\subsection{Sources}
\begin{itemize}
\item Local UnsolvedMath registry, supplied complete record: permanent ID \(30005582\), title \emph{Ricci Flow Uniqueness from Metric Space Initial Data}, exact source problem code \texttt{catalogue entry}; extracted passage labeled \emph{Problem 1}. The record belongs to the Oberwolfach Reports Open Problems collection. No individual proposer is supplied.
\item Supplied report citation: \emph{Differentialgeometrie im Grossen (2024)}, Oberwolfach Reports, pp.\ 1617--1670, DOI \texttt{10.4171/owr/2023/29}, \url{https://doi.org/10.4171/owr/2023/29}. The supplied background describes Oberwolfach Reports as documenting mathematical workshops and their open-problem discussions. The record gives proposed year \(2023\); the citation's year \(2024\) and the DOI's \(2023\) component are retained without reconciliation.
\item Supplied literature source, labeled Deruelle--Schulze--Simon: \url{https://wrap.warwick.ac.uk/id/eprint/190769/}. No paper title, theorem number, or complete bibliographic citation is supplied for this link.
\item Supplied literature source, labeled original OWR report: \url{https://doi.org/10.4171/owr/2023/29}.
\item Literature assessment dated \texttt{2026-08-22}: this is supplied dated curation, not independently verified. The record labels the problem \texttt{partially\_solved} and reports that stability results give effective comparison for Ricci flows from Reifenberg and locally bi-Lipschitz initial spaces under additional curvature hypotheses, while the broader regularity-and-uniqueness problem remains open. The supplied background expressly requests independent verification of time-sensitive status and cited claims. The additional curvature hypotheses are not stated in the record.
\item The extracted passage refers to its internal reference \([5,\text{ Theorem }1.2]\) for a counterexample concerning H\"older or Lipschitz Riemannian realization of the limiting distance. Its bibliographic identity is not supplied.
\item The registry marks the statement \texttt{corrected\_verified} and records that the official report was checked directly. This is a supplied source-verification claim; no independent source visit or verification is asserted in this draft.
\end{itemize}
\end{document}
